\documentclass[12pt, a4paper]{article}

\usepackage[utf8]{inputenc}
\usepackage[T1]{fontenc}
\usepackage[spanish]{babel}
\usepackage{geometry}
\usepackage{booktabs}
\usepackage{array}
\usepackage{longtable}
\usepackage{hyperref}
\usepackage{titlesec}
\usepackage{parskip}

\geometry{margin=2.5cm}

\hypersetup{
    colorlinks=true,
    linkcolor=blue,
    urlcolor=blue
}

\title{\textbf{Ruta de Aprendizaje}\\[0.5em]
       \large IA Aplicada a la Industria}
\author{Mg. Myriam Vera Fuentes}
\date{Duración: 16 semanas}

\begin{document}

\maketitle
\thispagestyle{empty}

\section{Metodología}

El curso se estructura en clases teóricas presenciales, centradas en la exposición de conceptos clave, análisis de casos reales y discusión crítica sobre el impacto y la implementación de la Inteligencia Artificial en la industria.

Adicionalmente, los estudiantes desarrollarán laboratorios de manera independiente, los cuales consisten en la ejecución de programas en Python entregados por la docente. Estos programas ya están implementados, por lo que no se requiere que los estudiantes programen.

La tarea del estudiante será ejecutar estos programas en el entorno provisto (Google Colab o Jupyter Notebook) y realizar una interpretación técnica y estratégica de los resultados obtenidos, considerando su aplicación en contextos industriales reales.

El curso culmina con un \textbf{proyecto final}, enfocado en la comprensión e interpretación de una solución basada en IA aplicada a un caso del mundo industrial.

\section{Evaluación}

\begin{itemize}
    \item \textbf{Informes de Laboratorios:} 50\% \quad (trabajo individual o grupal).
    \item \textbf{Proyecto Final:} 50\% \quad (análisis e interpretación de una solución IA industrial).
\end{itemize}

\section{Ruta de Aprendizaje Semanal}

\begin{longtable}{>{\centering\arraybackslash}p{1cm} >{\centering\arraybackslash}p{1.2cm} p{3.2cm} p{6cm} p{3cm}}
\toprule
\textbf{Sem.} & \textbf{Unidad} & \textbf{Tema} & \textbf{Contenido Clave} & \textbf{Evaluación} \\
\midrule
\endfirsthead
\toprule
\textbf{Sem.} & \textbf{Unidad} & \textbf{Tema} & \textbf{Contenido Clave} & \textbf{Evaluación} \\
\midrule
\endhead
\bottomrule
\endfoot

1  & U1 & Introducción a la IA
   & Historia, evolución, tipos de IA, IA vs.\ ML vs.\ DL, casos industriales
   & --- \\[4pt]

2  & U2 & Pensamiento estratégico en IA
   & IA como ventaja competitiva, impacto en decisiones y eficiencia
   & --- \\[4pt]

3  & U3 & Fundamentos de los datos
   & Recolección, limpieza, normalización, calidad de datos
   & Informe Lab 1 \\[4pt]

4  & U3 & Datos para entrenamiento
   & Representación adecuada, datos tabulares, no estructurados
   & Informe Lab 2 \\[4pt]

5  & U4 & Aprendizaje automático
   & Supervisado vs.\ no supervisado. Clasificación y regresión
   & Informe Lab 3 \\[4pt]

6  & U4 & Aplicaciones industriales del ML
   & Casos: predicción de demanda, control de calidad, logística
   & Informe Lab 4 \\[4pt]

7  & U5 & El perceptrón
   & Modelo de neurona, funciones de activación, aprendizaje básico
   & Informe Lab 5 \\[4pt]

8  & U5 & Redes neuronales multicapa
   & Arquitecturas MLP, hiperparámetros, concepto de caja negra
   & Informe Lab 6 \\[4pt]

9  & U6 & Arquitecturas avanzadas
   & CNN, RNN y Transformers. Casos: visión artificial, lenguaje, mantenimiento
   & Informe Lab 7 \\[4pt]

10 & U6 & Transferencia y escalamiento
   & Transfer learning, adaptabilidad de modelos, casos reales
   & Informe Lab 8 \\[4pt]

11 & U7 & Evaluación de modelos IA
   & Validación cruzada, métricas, ROI y KPIs industriales
   & Informe Lab 9 \\[4pt]

12 & U7 & Interpretabilidad y confianza
   & Modelos explicables, ética de la decisión automatizada
   & Informe Lab 10 \\[4pt]

13 & U8 & Aspectos éticos y regulatorios
   & Sesgos, privacidad, normativas y gobernanza en IA
   & Inicio Proyecto Final \\[4pt]

14 & U8 & IA Responsable en la industria
   & Responsabilidad corporativa, sostenibilidad y transparencia
   & Desarrollo Proyecto \\[4pt]

15 & U9 & Oportunidades emergentes
   & IA generativa, escalabilidad, IA industrial 4.0
   & Proyecto en curso \\[4pt]

16 & U10 & Presentación final
   & Presentación de resultados interpretados y conclusiones estratégicas
   & \textbf{Entrega Proyecto Final (50\%)} \\

\end{longtable}

\end{document}